\section*{Methods}

% ─────────────────────────────────────────────────────────────────────────────
\subsection*{Participants and ethics}
% ─────────────────────────────────────────────────────────────────────────────

Data were collected as part of a previously published study~\cite{espinoza2025}.
Thirty-six healthy adults (18--35~years; natural dentition; no dental pain,
temporomandibular disorder diagnosis, or medications affecting neuromuscular
function) participated.
All procedures were approved by the Ethics Committee of the School of Dentistry,
Universidad de Valpara\'{i}so (registration code POSTG-06-22) and conducted in
accordance with the Declaration of Helsinki.
Written informed consent was obtained from all participants before inclusion;
the present secondary analysis falls within the scope of that approval.

% ─────────────────────────────────────────────────────────────────────────────
\subsection*{Study design and anesthesia protocol}
% ─────────────────────────────────────────────────────────────────────────────

The parent study used a 2\,$\times$\,2 within-subject crossover design crossing
chewing (chewing vs.\ no-chewing) with anesthesia condition (anesthesia vs.\
placebo), counterbalanced in two blocks (R1: anesthesia block first; R2: placebo
block first).
The present analysis focuses exclusively on the chewing blocks.
At the start of each chewing block, the experimenter applied one dose of either
topical lidocaine (10\% pump spray) or a matched placebo spray (identical
excipients and taste, no active agent) to the mucosa and periodontal tissues
ipsilateral to the non-chewing side.
No reapplication occurred during the block.
Participants then completed four 60-second chewing bouts with a standardised
food item, separated by approximately 4.5-minute rest intervals, yielding bout
onset times of approximately 3, 7.5, 12, and 17~minutes post-application
per condition.

% ─────────────────────────────────────────────────────────────────────────────
\subsection*{EMG recording}
% ─────────────────────────────────────────────────────────────────────────────

Electrophysiological data were recorded continuously with a BioSemi ActiveTwo
system (BioSemi, Amsterdam, Netherlands) at 1024~Hz from 64 scalp electrodes
and 8 external electrodes (ExG1--8).
Masseter EMG was captured from ExG5--8 using Ag/AgCl electrodes positioned over
the masseter belly in a bipolar configuration following established surface EMG
methodology~\cite{deLuca1997}:
left bipolar pair (ExG5\,$-$\,ExG6) and right bipolar pair (ExG7\,$-$\,ExG8).

% ─────────────────────────────────────────────────────────────────────────────
\subsection*{Signal processing}
% ─────────────────────────────────────────────────────────────────────────────

EMG processing was implemented in Python~3.10 using NumPy, SciPy~\cite{SciPy2020},
and custom modules (analysis code will be deposited at a public repository
upon acceptance; available from the corresponding author on request).
Raw BDF files were processed at the original 1024~Hz sampling rate without
resampling; resampling was avoided to preserve spectral content up to 450~Hz and
prevent aliasing artefacts in the median frequency estimate.
The bipolar derivations were computed as described above.
Signals were band-pass filtered between 20 and 450~Hz using a zero-phase
4th-order Butterworth filter.
Line-frequency noise and harmonics (50, 100, 150, and 200~Hz) were removed with
the Zapline-plus algorithm~\cite{Pyzaplineplus}, which attenuates narrowband
interference without distorting the spectrum at adjacent frequencies, unlike
conventional notch filters.
For each participant, the recording side (left or right) with EMG exhibiting
physiological median frequency (55--170~Hz) and chewing rate (0.6--2.6~Hz)
in both conditions, and with the higher signal-to-noise ratio, was selected for
analysis.
Two participants (M3, M7) were excluded for confirmed dead electrode
artefact in one condition.

% ─────────────────────────────────────────────────────────────────────────────
\subsection*{Cycle detection and feature extraction}
% ─────────────────────────────────────────────────────────────────────────────

Chewing bouts (60~s duration; onset at event code~1) were segmented from the
continuous recording.
Individual masticatory cycles were detected by computing a 100-ms RMS envelope
of the full-wave-rectified bipolar signal, normalising it to a rest-period
$z$-score ($z = (\text{env} - \mu_{\text{rest}}) / \sigma_{\text{rest}}$), and
identifying peaks in the $z$-score trace above the within-bout 50th percentile
with a minimum inter-peak distance of 0.4~s (SciPy~\cite{SciPy2020}).
Using a rest-normalised relative threshold rather than a fixed amplitude cutoff
ensures equivalent detection sensitivity across conditions regardless of signal
magnitude.
To verify that cycle detection was not differentially biased by the amplitude
difference between conditions, we confirmed that the mean number of cycles
detected per bout was equivalent between anesthesia and placebo at every bout
(bout~1: $\bar{n}$\,=\,87.8 vs.\ 87.1 cycles; bouts~2--4: differences
$\leq$\,3 cycles; no paired test approached significance).

Per-bout, per-participant features were extracted: median cycle amplitude
(\textit{med\_amp}), defined as the within-bout median of the instantaneous
amplitude of the full-wave-rectified signal at each detected cycle peak time;
RMS amplitude, total signal power, power in 20--60~Hz and
60--150~Hz bands (Welch method), median frequency (MDF) and mean frequency (MNF)
from the full power spectral density, chewing rate (cycles per second),
coefficient of variation of the inter-peak interval (CV\textsubscript{IPI}), and
duty cycle.
Cycle-locked time-frequency maps were computed by epoching the bipolar signal at
$[-0.4, +0.4]$~s around each detected peak and applying Morlet wavelet
convolution at 30 logarithmically spaced frequencies from 20 to 450~Hz
(wavelet cycles\,=\,7).

% ─────────────────────────────────────────────────────────────────────────────
\subsection*{Statistical analysis}
% ─────────────────────────────────────────────────────────────────────────────

All inferential tests used $N$\,=\,34 participants.
Within each chewing bout, the anesthesia-placebo difference for each metric was
tested with a two-sided paired Wilcoxon signed-rank test~\cite{Pingouin}.
Multiple comparisons across 11 metrics within each bout were controlled by the
Benjamini--Hochberg false-discovery rate (FDR) procedure~\cite{Benjamini1995}.
FDR correction was applied within each bout separately, consistent with the
a~priori hypothesis that bout~1 would carry the primary pharmacological effect;
the four bout comparisons were treated as planned time-points (not as a joint
family requiring cross-bout correction).
Effect sizes are reported as paired Cohen's $d_z$~\cite{Cohen1988}.
A post-hoc sensitivity analysis (paired $t$ equivalent, $\alpha$\,=\,0.05
two-tailed, $1-\beta$\,=\,0.80, $N$\,=\,34) indicated that the study was
powered to detect effects of $d_z$\,$\geq$\,0.49; the primary bout~1 effect
($d_z$\,=\,0.695) substantially exceeds this threshold, while the null effects
at bouts~3--4 ($d_z$\,$\leq$\,0.10) are well below it, consistent with the
TOST-confirmed equivalence at those bouts.
Session-averaged equivalence of timing metrics was assessed with two one-sided
tests (TOST) at an equivalence bound of $|d|$\,=\,0.5~\cite{Lakens2017};
this bound corresponds to a medium effect size and represents a conventional
threshold for practical equivalence that is accepted as meaningful in the
sensorimotor literature~\cite{Lakens2017}.
The condition\,$\times$\,bout interaction was tested with a repeated-measures
ANOVA using Pingouin~\cite{Pingouin} with Greenhouse--Geisser sphericity
correction.
Cycle-locked TF maps were compared between conditions by a one-sample
cluster-based permutation test on the paired ANE\,$-$\,PLA difference map
(5\,000 permutations; cluster-forming threshold $t$ at $p$\,=\,0.05)~\cite{Maris2007}.
Crossover confound analyses used a mixed ANOVA (condition\,$\times$\,counterbalancing
arm) at bout~1 and an unpaired Mann--Whitney test for residual carryover.
As an exploratory analysis, bilateral masseter synchrony was quantified for
participants with physiologically valid signals on both left and right sides
($N$\,=\,31; median frequency 45--200~Hz and chewing rate 0.4--2.8~Hz in
both conditions).
Within each bout, the Pearson correlation between left and right masseter
RMS envelopes (50-ms window) was computed as a measure of bilateral temporal
coupling, and the amplitude asymmetry index (L\,$-$\,R)/(L\,$+$\,R) was
derived from the mean envelope amplitudes.
Bilateral metrics were compared between conditions using paired Wilcoxon tests
(FDR-corrected) and a 2\,$\times$\,4 repeated-measures ANOVA
(condition\,$\times$\,bout).

